\documentclass[11pt]{article}
\usepackage[margin=1in]{geometry}
\usepackage{amsmath,amssymb,amsthm,mathtools,booktabs,hyperref,natbib}
\newtheorem{lemma}{Lemma}
\newtheorem{proposition}{Proposition}
\newtheorem{corollary}{Corollary}
\newtheorem{remark}{Remark}
\title{Regional Public Breeding and Variety Replacement: Decentralized Research Incentives and National Welfare}
\author{Anonymous draft for internal research verification}
\date{October 2026 --- Stage 10 working manuscript}
\begin{document}
\maketitle
\begin{abstract}
Regional governments may fund the development of improved crop varieties while growers independently decide whether to replace established cultivars. This paper studies a two-stage game in which territorial public laboratories choose quality-improving research efforts and competitive farmers subsequently allocate fixed land between old and new varieties. We derive the adoption response and characterize the discrepancy between decentralized public research and national welfare. The research wedge changes sign with equilibrium variety adoption, even when the growers' conditional allocation is efficient. In an explicit parameter family, each regional government has a strictly positive, globally optimal research response while the unique national optimum entails no research. The resulting decentralized equilibrium lowers national welfare relative to no research. We establish this result across all planting regimes, identify its exact parameter boundary, and delineate the demand-curvature and institutional restrictions under which the mechanism generalizes. The analysis is theoretical; it does not estimate actual regional breeding returns.
\end{abstract}

\noindent\textbf{Keywords:} public agricultural research; variety replacement; regional government competition; endogenous adoption; national welfare.\\
\textbf{JEL:} H41, H77, O32, Q16.

\section{Introduction}
Public breeding programs are often organized territorially, while agricultural product markets extend beyond administrative boundaries. A new cultivar can improve its region's marketed product, yet its commercial success may displace returns earned by producers elsewhere. The resulting fiscal incentive is not captured by comparing technical research productivity alone: public research affects market outcomes through the subsequent planting decisions of competitive growers.

This paper asks when individually rational regional public breeding decisions diverge from the nationally efficient research allocation after growers endogenously replace established varieties. The key distinction is between the quality-investment margin, controlled by public laboratories, and the adoption margin, controlled by growers. A region's producer rents and its assigned share of consumer surplus need not internalize the same price effects as a national planner.

The first result establishes that, conditional on research quality, competitive planting can be efficient even though the preceding regional research decisions are not. The second result derives a sign-changing regional--national research wedge. The third establishes a stronger possibility: in a transparent analytically solvable economy, positive research is the unique decentralized equilibrium, but zero research is the unique national optimum. The exact global proof includes zero, partial and complete adoption, rather than relying on interior first-order conditions.

These are claims about a specified game, not a claim that excessive public R\&D, rent shifting, or beggar-thy-neighbor competition is new in itself. Private-profitability and welfare-divergence results already arise in the advertising literature \citep{dixit1978advertising,alston2001beggar,cardon2003market}; related markets and competitive resource reallocation have also been studied \citep{kinnucan2000distributional,balagtas2007dairy,lence2007collective}. In particular, a socially optimal zero level of a private promotion instrument is not, by itself, novel. The potential contribution is the explicitly solved interaction of territorial public quality provision, constrained competitive varietal replacement, and the full national-welfare benchmark, including global boundary cases. The existing source-level assessment does not yet rule out absorption by all relevant parent theorems.

\section{Economic environment}
\subsection{Timing, goods, and public objectives}
There are two regions, $i\in\{1,2\}$, each with one unit of agricultural land. The old and new varieties in region $i$ are denoted $O_i$ and $N_i$. In stage 1, the regional public laboratory chooses $x_i\geq 0$, a quality-enhancing research effort for $N_i$. In stage 2, price-taking growers choose the share $a_i\in[0,1]$ of land planted to $N_i$, so quantities are $q_{O_i}=1-a_i$ and $q_{N_i}=a_i$. Consumer markets then clear. Governments correctly anticipate the competitive stage-2 outcome.

For $q=(q_{O_1},q_{N_1},q_{O_2},q_{N_2})^\top$, national gross utility is
\begin{equation}
 V(q;x)=\alpha(x)^\top q-\tfrac12q^\top Bq,\quad
 \alpha(x)=(A,A+\Delta+x_1,A,A+\Delta+x_2)^\top,
 \label{eq:utility}
\end{equation}
where
\begin{equation}
 B=\begin{pmatrix}
 1&t&b&c\\ t&1&c&n\\ b&c&1&t\\ c&n&t&1
 \end{pmatrix},\qquad p=\nabla_q V=\alpha(x)-Bq.
 \label{eq:demand}
\end{equation}
Assume $B$ is positive definite in the relevant parameter domain and all realized prices are economically admissible. Regional growers bear real conversion cost $K_i(a_i)=k_0a_i+\kappa a_i^2/2$ and the regional government bears $C_i(x_i)=Rx_i^2/2$, with $R,\kappa>0$. Let $\Pi_i$ be regional producer receipts net of $K_i$, and let $CS=V-p^\top q$ be aggregate consumer surplus. Region $i$ maximizes
\begin{equation}
 G_i=\Pi_i+\tfrac12 CS-C_i(x_i),
 \label{eq:local}
\end{equation}
whereas national welfare is
\begin{equation}
 W=\Pi_1+\Pi_2+CS-C_1(x_1)-C_2(x_2)=G_1+G_2.
 \label{eq:welfare}
\end{equation}
The equal division of consumer surplus is a transparent accounting assumption, not an empirical description of regional incidence.

\subsection{Competitive adoption}
Define
\begin{equation}
 h=1-t,\quad j=b-c,\quad e=b+n-2c,\quad
 d=\kappa+2h,\quad M=d+e,\quad L=d-e,\quad
 z=\Delta+h+j-k_0.
 \label{eq:notation}
\end{equation}
The new-versus-old producer price advantage net of marginal conversion cost is
\begin{equation}
 p_{N_i}-p_{O_i}-K_i'(a_i)=z+x_i-d a_i-e a_j,\quad j\ne i.
 \label{eq:adoptionfoc}
\end{equation}
The stage-2 equilibrium obeys the complementary conditions: this residual is nonpositive at $a_i=0$, zero for $0<a_i<1$, and nonnegative at $a_i=1$. For an interior allocation,
\begin{equation}
 \begin{pmatrix}d&e\\e&d\end{pmatrix}
 \begin{pmatrix}a_1\\a_2\end{pmatrix}
 =\begin{pmatrix}z+x_1\\z+x_2\end{pmatrix}.
 \label{eq:adoptionmatrix}
\end{equation}
The positive-definiteness conditions $M,L>0$ imply uniqueness of the conditional adoption equilibrium.

\begin{lemma}[Conditional efficiency of competitive adoption]
Fix research efforts $(x_1,x_2)$. If the quadratic surplus objective is strictly concave in the planting shares and growers face the same real conversion costs as the planner, the competitive stage-2 allocation globally maximizes national welfare over $(a_1,a_2)\in[0,1]^2$ conditional on $x$.
\end{lemma}
\begin{proof}
Differentiating $V(q(a);x)-\sum_iK_i(a_i)$ with respect to $a_i$ gives $p_{N_i}-p_{O_i}-K_i'(a_i)$. Thus the planner's constrained first-order conditions coincide with the competitive growers' KKT conditions, including both planting corners. Strict concavity makes these conditions globally sufficient. The quadratic form and KKT implication are independently formalized in \texttt{Core.lean}.
\end{proof}
The lemma does not establish efficient stage-1 research. It instead isolates the externality in governments' incentives to shift the subsequent market allocation.

\section{Decentralized research incentives}
For a symmetric interior allocation, $a_1=a_2=a=(z+x)/M$. Define
\begin{equation}
 \xi=\frac{2h-e}{2L},\qquad a^\star=\frac{h-j}{2h-e},
 \label{eq:threshold}
\end{equation}
assuming the denominators are nonzero and the positive-sign restrictions stated below hold.

\begin{proposition}[Regional--national research wedge]
At a symmetric interior equilibrium of the planting subgame, the difference between the regional government's marginal objective and the national planner's marginal welfare with respect to $x_i$ is
\begin{equation}
 \frac{dG_i}{dx_i}-\frac{dW}{dx_i}=\xi(a^\star-a).
 \label{eq:wedge}
\end{equation}
When $\xi>0$ and $a^\star\in(0,1)$, the wedge is positive for $a<a^\star$ and negative for $a>a^\star$. An interior threshold is essential for a sign change within the feasible adoption domain. The formula is an interior, linear-demand statement; it does not assert that the global research equilibrium is interior for every parameter value.
\end{proposition}
\begin{proof}
To make the economic-to-algebra bridge explicit, competitive prices, producer income and consumer surplus on the \emph{interior follower branch} can be reduced to
\begin{align}
p_{iO}&=A-1-b+h a_i+j a_j,\nonumber\\
p_{iN}&=A+\Delta+x_i-t-c-h a_i-(n-c)a_j,\nonumber\\
\Pi_i&=A-1-b+h a_i+j a_j+\frac{\kappa}{2}a_i^2,\nonumber\\
CS&=1+b-(h+j)(a_1+a_2)
+h(a_1^2+a_2^2)+e a_1a_2.
\label{eq:incidencebridge}
\end{align}
The producer-income simplification invokes the interior grower FOC; none of these interior-reduced expressions is used to assess off-equilibrium planting corners in Proposition 3.

The inverse adoption matrix in \eqref{eq:adoptionmatrix} gives $\partial a_i/\partial x_i=d/(d^2-e^2)$ and $\partial a_j/\partial x_i=-e/(d^2-e^2)$. The national welfare envelope derivative is $dW/dx_i=a_i-Rx_i$, because adoption maximizes national welfare conditional on $x$. Differentiate \eqref{eq:local} with \eqref{eq:incidencebridge} and these adoption responses. At symmetry, the matrix eigenvalues are $M$ and $L$. Collecting the local-minus-national marginal effect yields
\[
 \frac{h-j-(2h-e)a}{2L}
 =\frac{2h-e}{2L}\left(\frac{h-j}{2h-e}-a\right).
\]
This is \eqref{eq:wedge}. The algebraic factorization is kernel-checked in \texttt{Core.lean}; the payoff-incidence derivation remains visible here.
\end{proof}

Assume $RM>1$, $RM-1+\xi>0$, and the following candidates are feasible interior solutions and are globally optimal in their respective stage-1 problems. The symmetric national planner and decentralized research candidates are
\begin{equation}
 x^P=\frac{z}{RM-1},\qquad
 x^D=\frac{(1-\xi)z+\xi M a^\star}{RM-1+\xi}.
 \label{eq:efforts}
\end{equation}
The expression for $x^P$ is applicable only where its interior solution is feasible; in particular, it must not be used to claim a negative optimal effort when $z<0$.

\begin{corollary}[Research and adoption gaps]
Under the preceding interior and globality conditions,
\begin{equation}
 x^D-x^P=\frac{\xi M}{RM-1+\xi}(a^\star-a^P),
 \qquad a^D-a^P=\frac{x^D-x^P}{M}.
 \label{eq:corollary}
\end{equation}
\end{corollary}
\begin{proof}
Subtract the two research first-order conditions in \eqref{eq:efforts} and use $a^P=(z+x^P)/M=Rx^P$. The second equality follows directly from the symmetric adoption equation. The corresponding algebra is checked in \texttt{Core.lean}.
\end{proof}

When both benchmarks are indeed global interior solutions on the same symmetric adoption branch, the national-welfare cost has the additional closed form
\begin{equation}
W(x^P,x^P)-W(x^D,x^D)
=\left(R-\frac{1}{M}\right)(x^D-x^P)^2\geq0,
\label{eq:interiorwelfarecost}
\end{equation}
because $W(x,x)$ is a concave quadratic with vertex $x^P$ under $RM>1$. The claim is a \emph{conditional equilibrium comparison}, not an extension to the corner planner in Proposition 3. This square-completion identity is separately checked in \texttt{Core.lean}.

\subsection{Demand curvature and the scope of the sign result}
The nonlinear result concerns a separate, explicitly restricted class with $N\geq2$ identical regions, unit land per region, and a national quasilinear gross-utility function
\begin{equation}
V(q;x)=\sum_{i=1}^{N}\sum_{s\in\{O,N\}}\alpha_{is}(x_i)q_{is}
-\sum_{i=1}^{N}\sum_{s\in\{O,N\}}\phi(q_{is})
-\frac{\beta}{2}Q^2-\frac{r}{2}\sum_iQ_i^2
-\frac{g}{2}(Q_O^2+Q_N^2).
\label{eq:nonlinearutility}
\end{equation}
Here $Q_i=q_{iO}+q_{iN}=1$, $Q=\sum_iQ_i=N$, $Q_s=\sum_iq_{is}$, $\beta,r,g\geq0$, and $\phi$ is twice continuously differentiable with $\phi''>0$. Only $\alpha_{iN}$ receives the own-region quality increment $f(x_i)$, with $f'>0$; old and new varieties must share \emph{the same} $\phi$. Assume $K$ is convex and twice differentiable, all planting choices are interior and symmetric, and each regional authority receives $CS/N$. These restrictions are part of the proposition, not innocuous normalization.

Write $H_\phi(q)=q\phi''(q)$ and $L_\phi=K''(a)+\phi''(1-a)+\phi''(a)>0$. The planting Jacobian is $L_\phi I_N+2g\mathbf1\mathbf1^\top$; its own-minus-average response to one region's research shock is $f'(x)(N-1)/(NL_\phi)$. The precise local--national research wedge is
\begin{equation}
G_{i,x_i}-W_{x_i}
=\frac{f'(x)(N-1)}{NL_\phi}
\bigl[H_\phi(1-a)-H_\phi(a)\bigr].
\label{eq:nonlinearwedge}
\end{equation}
Under $f'(x)>0$ and $L_\phi>0$, the sign therefore equals the sign of
\begin{equation}
 H(1-a)-H(a),
 \label{eq:curvature}
\end{equation}
provided the remaining prefactors are positive.

\begin{proposition}[Restricted nonlinear-demand incidence]
For the symmetric interior adoption equilibrium of the restricted anonymous separable-demand class, and assuming an invertible adoption Jacobian and positive response multipliers, the curvature contribution to the local--national research wedge changes sign according to \eqref{eq:curvature}. In particular, adoption share $a=1/2$ is curvature-neutral when old and new own-demand curvature functions coincide.
\end{proposition}
\begin{proof}
Differentiate the competitive farmer planting equations implied by \eqref{eq:nonlinearutility}. Their Jacobian in the adoption shares is $L_\phi I_N+2g\mathbf1\mathbf1^\top$, whose eigenvalue on the zero-sum planting subspace is $L_\phi$. Subtracting any nonshocked-region response from the shocked-region response therefore cancels the common $g$ component. Averaging across regions gives the exact own-minus-average response $f'(x)(N-1)/(NL_\phi)$. Differentiating each region's producer receipts and its $1/N$ consumer-surplus share, while applying the national welfare envelope theorem, leaves precisely the price-incidence term $H_\phi(1-a)-H_\phi(a)$ multiplied by that response. This yields \eqref{eq:nonlinearwedge} and \eqref{eq:curvature}. \texttt{Prop2N.lean} verifies the conditional response algebra in Lean; the economic utility-to-Jacobian derivation is provided here, not claimed as a separate Lean proof.
\end{proof}
Asymmetric old/new curvature can produce a nonzero wedge even at half adoption. Thus this proposition is not a universal theorem for arbitrary nonlinear demand.

\section{When regional breeding research lowers national welfare}
The interior wedge alone does not establish globally excessive research. This section supplies an exact economy in which all adoption corners and large unilateral research deviations are controlled.

Set
\begin{equation}
 t=\tfrac12,\quad b=c=n=\tfrac15,\quad
 k_0=\kappa=1,\quad R=2,\quad A=10,\quad
 \Delta=\tfrac12+z.
 \label{eq:family}
\end{equation}
Then $e=j=0$, and the competitive planting share in each region is
\begin{equation}
 a_i(x_i)=\min\left\{1,\max\left\{0,\frac{z+x_i}{2}\right\}\right\}.
 \label{eq:clamp}
\end{equation}
In this family, the regions' first-stage best responses are independent of the other region's effort. This is a legitimate special case, but it does not demonstrate strict strategic substitutability between public research efforts.

Let $z_c=(-4+\sqrt{13})/12$.

\begin{proposition}[Globally harmful decentralized public breeding]
For every $z\in(z_c,0]$ in \eqref{eq:family}, the unique regional research Nash equilibrium is symmetric with
\begin{equation}
 x^D=\frac{1+3z}{13}>0,\qquad
 a^D=\frac{1+16z}{26}>0.
 \label{eq:p3nash}
\end{equation}
The unique national optimum has $x_1=x_2=a_1=a_2=0$. The welfare difference is strictly negative:
\begin{equation}
 W^D-W^P=\frac{(10z-1)(22z+3)}{338}<0.
 \label{eq:p3loss}
\end{equation}
At $z=z_c$ the regional best-response set becomes $\{0,\sqrt{13}/52\}$, yielding four pure research Nash profiles; hence the lower inequality is strict.
\end{proposition}

\begin{proof}
We derive the regional payoff from \emph{unreduced} utility and prices, not from an interior follower condition extrapolated to corners. Relative to the zero-research, zero-adoption baseline, a regional government's payoff at arbitrary planting share $a_i$ is
\begin{equation}
 \mathcal F_i(x_i,a_i)
 =(z+x_i+\tfrac14)a_i-\tfrac54a_i^2-x_i^2.
 \label{eq:offpathpayoff}
\end{equation}
Substituting \eqref{eq:clamp} gives three regimes:
\begin{equation}
 F_i(x_i)=
 \begin{cases}
 -x_i^2,&0\leq x_i\leq-z,\\[2pt]
 (z+x_i)/8+3(z+x_i)^2/16-x_i^2,&-z<x_i<2-z,\\[2pt]
 z+x_i-1-x_i^2,&x_i\geq2-z.
 \end{cases}
 \label{eq:threebranches}
\end{equation}
The middle branch is strictly concave, with unique stationary maximizer $x^D=(1+3z)/13$ and payoff
\[
 F_i(x^D)=\frac{48z^2+32z+1}{208}>0
 \quad\text{for }z>z_c.
\]
The zero-adoption branch never exceeds zero. The complete-adoption branch is decreasing over $x_i\geq2-z$ because its derivative is $1-2x_i<0$. Its boundary value is strictly below the middle-branch maximizer. Hence \eqref{eq:p3nash} is the unique \emph{global} regional best response, and the two independent responses form the unique research Nash equilibrium.

For national welfare, without imposing either planting equilibrium or symmetry, the difference from the zero allocation is
\begin{align}
 W-W_0
 &=z(a_1+a_2)+x_1a_1+x_2a_2
   -(a_1^2+a_2^2)-(x_1^2+x_2^2)\nonumber\\
 &=z(a_1+a_2)-\frac34(a_1^2+a_2^2)
   -(x_1-a_1/2)^2-(x_2-a_2/2)^2
 \leq0.
 \label{eq:plannerbound}
\end{align}
For $z\leq0$ and $a_i\geq0$, equality requires $a_i=x_i=0$ in both regions. This proves unique national optimality over a feasible set \emph{larger} than the competitive follower allocations. Substituting \eqref{eq:p3nash} into \eqref{eq:plannerbound} yields \eqref{eq:p3loss}. At the excluded endpoint, $F_i(0)=F_i(\sqrt{13}/52)=0$, and no other nonnegative effort is optimal. The proof-critical three-regime inequalities, nationwide square completion, uniqueness and cutoff are formally checked in \texttt{Prop3Global.lean}, \texttt{Uniqueness.lean}, and \texttt{Boundary.lean}; \texttt{PrimitivesBridge.lean} checks the exact off-equilibrium utility-to-payoff reduction.
\end{proof}

The result is stronger than a local overinvestment derivative: each region finds positive research privately optimal, while the national planner strictly prefers no research. Nevertheless, it is an \emph{existence theorem} for an explicit admissible family, not a statement that all regional breeding systems overinvest.

\section{Robustness, limitations, and institutional interpretation}
\subsection{Meaningful alternative specifications}
The pre-freeze falsification exercises show that the harmful-research phenomenon survives a convex quartic public research cost in a specified re-solved alternative, while the half-adoption neutrality in the nonlinear demand proposition fails under asymmetric old/new own-demand curvature. A separate quartic conversion-cost exercise retains the sign of the regional--national wedge at its tested interior allocation. These exercises support conditional portability; they do not imply general robustness across all cost and demand classes.

\subsection{Nonzero cross-region interaction}
The analytic family in Section 4 has $e=0$. A distinct rational parameter example with $e=1/100$ yields positive symmetric public research $x^D=9751/130096$ and adoption $a^D=1051/32524$, with the national optimum again at zero and a negative welfare difference $-1933000423/211562115200$. Its local cross-research derivative is $-50/39999<0$. The global unilateral best response was checked by exact arithmetic across the nine adoption KKT regimes, whereas the Lean module \texttt{ExactExamples.lean} certifies the exact candidate and planner-welfare identities, \emph{not} a full Lean proof of nine-regime Nash globality or uniqueness of all asymmetric equilibria.

A distinct, earlier benchmark uses $A=10$, $t=1/2$, $b=n=3/5$, $c=9/20$, $k_0=\kappa=1$, $R=2$, and $\Delta=17/50$ (so $e=3/10$). Exact piecewise unilateral-deviation verification gives the symmetric Nash research effort $x^D=389/6470$ and planting share $a^D=141/6470$; the national optimum is again zero, with welfare loss $2569157/418609000$. The example has nonzero competitive coupling and is not the analytic $e=0$ family of Proposition 3. Here too, Lean checks rational candidate and planner identities but not the complete piecewise Nash certificate.

\subsection{Interpretation and external validity}
A regional public breeder may rationally finance a quality improvement because the induced private variety replacement redistributes agricultural product-market rents. The national welfare effect also counts displaced rents and the real fiscal cost of research. The analysis does not imply that actual public breeding budgets should be eliminated: it abstracts from disease resistance, climate adaptation, spillovers to future varieties, uncertain discoveries, farmer heterogeneity, and institutional intergovernmental transfers. The model is a diagnostic benchmark for distinguishing adoption efficiency from research efficiency.

\section{Discussion}
The regional--national divergence does not require a market failure in the growers' conditional planting choices. It emerges when public laboratories anticipate an efficient competitive adoption response but assess its price and surplus incidence territorially. This separation clarifies why evidence that farmers efficiently adopt superior varieties cannot by itself establish that decentralized public research spending is efficient.

A natural policy question is whether interregional transfers or coordination could internalize the displaced surplus. Such mechanisms are not solved in the frozen model and are not claimed as results of this paper. Any policy prescription would additionally require empirical estimates of the cross-regional substitution parameters, the opportunity cost of public research and the relevant incidence of consumer surplus.

\section{Conclusion}
This paper separates publicly chosen breeding quality from competitively chosen variety replacement. Competitive adoption can maximize national welfare conditional on research while decentralized regional laboratories have different ex ante incentives. In a tractable parameter family, positive research is the unique global decentralized equilibrium even though the unique national optimum is zero research. The proof covers planting corners and large deviations and identifies an exact boundary at which research equilibria become multiple. The conclusions are deliberately bounded by the model's demand, fiscal and equilibrium assumptions.

\appendix
\section{Derivation and proof crosswalk}
Equations \eqref{eq:utility}--\eqref{eq:adoptionmatrix} supply the primitive-to-follower bridge. Equations \eqref{eq:wedge}--\eqref{eq:corollary} are interior statements. Equations \eqref{eq:offpathpayoff}--\eqref{eq:plannerbound} supply the full global proof, including all adoption regimes. Formal artifacts are in the pinned Lean project \texttt{checkpoints/2026-10-24/18\_lean\_formal\_proofs/lean}; the last checked build and allowed axioms are documented in the Stage 8 freeze. The full paper-facing economic model, and in particular the generic nonlinear-demand-to-Jacobian bridge, is not claimed to be kernel-certified in its entirety.

\section{Further scope of the numerical checks}
For the interacting rational example, the existing exact-arithmetic program enumerates the nine follower KKT status combinations and tests unilateral research deviations across feasible branches. This establishes a checked candidate equilibrium under that audit, not a Lean proof of all asymmetric Nash equilibria. No new numerical results are introduced here.

\bibliographystyle{plainnat}
\bibliography{references}
\end{document}
